% GENERATED FILE. Edit canonical lessons, then run npm run generate:books.
% canonical-source-hash: fnv1a64:7f074f9d10a491a9
% canonical-lessons: LA-C545-sonus, LA-C545-sono, LA-C545-resono, LA-C545-symphonia, LA-C545-concentus

\chapter{Sound and harmony}
\label{ch:la-sound-and-harmony}
\hlchaptermodality{545}

\begin{chapteropening}
\textbf{By the end of this chapter:} \emph{I can describe sounds, echoes and harmony.}

Complete the last lesson of chapter 545: I can describe sounds, echoes and harmony.
\end{chapteropening}

\section[sonus, sonī]{sonus, sonī — a sound}

\label{lesson:LA-C545-sonus}



\begin{quote}
Before the new one: say the Latin for a lyre-player, then the Latin for to play the cithara.
\end{quote}

\subsection*{You'll want to know: sonus, sonī}
\textbf{sonus, sonī} — \textquotedblleft{}a sound\textquotedblright{}.

\begin{quote}
Sonus tubae per castra it.
\end{quote}

\begin{quote}
\emph{The sound of the trumpet goes through the camp.}
\end{quote}

Masculine, second declension. English \emph{sound} comes from it through French.

\begin{cousinweb}
\textbf{sonus} belongs with \emph{sonō} \textquotedblleft{}to sound\textquotedblright{}, inherited from a Proto-Indo-European root meaning \textquotedblleft{}to sound\textquotedblright{} (EtymDB; Lewis \& Short). English \emph{sound} comes from it through French.
\end{cousinweb}

\begin{grammarlens}[title={What you've built}]
One more word for telling, explaining and describing at B1.
\end{grammarlens}

\subsection*{Your turn}
\begin{itemize}
\raggedright
  \item \emph{Say it:} \emph{sonus, sonī}
  \item \emph{Say it:} \emph{sonus, sonī}, once more
  \item \emph{Say it:} a sentence of your own with \emph{sonus, sonī}
  \item \emph{Recall:} say the Latin for a lyre-player, then the Latin for to play the cithara, then say \emph{sonus, sonī} again
\end{itemize}

\begin{tcolorbox}[breakable,colback=teal!4,colframe=teal!35!black,title={Before you move on}]
What does sonus, sonī mean? (\textquotedblleft{}A sound\textquotedblright{}.) Say it once more.
\end{tcolorbox}
\section[sonō, sonāre]{sonō, sonāre — to sound, ring}

\label{lesson:LA-C545-sono}



\begin{quote}
Before the new one: say the Latin for to play the cithara, then the Latin for a sound.
\end{quote}

\subsection*{You'll want to know: sonō, sonāre}
\textbf{sonō, sonāre} — \textquotedblleft{}to sound, ring\textquotedblright{}.

\begin{quote}
Tintinnābulum sonat; hospitēs adsunt.
\end{quote}

\begin{quote}
\emph{The bell rings; the guests are here.}
\end{quote}

\emph{sonō, sonāre, sonuī, sonitum}. Italian \emph{sonata}, a piece to be played, comes from it.

\begin{cousinweb}
\textbf{sonō} is inherited from a Proto-Indo-European root meaning \textquotedblleft{}to sound\textquotedblright{} (EtymDB). Italian \emph{sonata} comes from it.
\end{cousinweb}

\begin{grammarlens}[title={What you've built}]
One more word for telling, explaining and describing at B1.
\end{grammarlens}

\subsection*{Your turn}
\begin{itemize}
\raggedright
  \item \emph{Say it:} \emph{sonō, sonāre}
  \item \emph{Say it:} \emph{sonō, sonāre}, once more
  \item \emph{Say it:} a sentence of your own with \emph{sonō, sonāre}
  \item \emph{Recall:} say the Latin for to play the cithara, then the Latin for a sound, then say \emph{sonō, sonāre} again
\end{itemize}

\begin{tcolorbox}[breakable,colback=teal!4,colframe=teal!35!black,title={Before you move on}]
What does sonō, sonāre mean? (\textquotedblleft{}To sound, ring\textquotedblright{}.) Say it once more.
\end{tcolorbox}
\section[resonō, resonāre]{resonō, resonāre — to resound, echo}

\label{lesson:LA-C545-resono}



\begin{quote}
Before the new one: say the Latin for a sound, then the Latin for to sound.
\end{quote}

\subsection*{You'll want to know: resonō, resonāre}
\textbf{resonō, resonāre} — \textquotedblleft{}to resound, echo\textquotedblright{}.

\begin{quote}
Tōta vallis cantū pāstōrum resonat.
\end{quote}

\begin{quote}
\emph{The whole valley echoes with the shepherds' song.}
\end{quote}

\emph{resonō, resonāre, resonāvī}: \emph{re-} 'back' + \emph{sonō}. English \emph{resound}, \emph{resonant}.

\begin{cousinweb}
\textbf{resonō} is built on \emph{re-} \textquotedblleft{}back\textquotedblright{} and \emph{sonō} \textquotedblleft{}to sound\textquotedblright{} (Lewis \& Short; EtymDB). English \emph{resound} and \emph{resonant} come from it.
\end{cousinweb}

\begin{grammarlens}[title={What you've built}]
One more word for telling, explaining and describing at B1.
\end{grammarlens}

\subsection*{Your turn}
\begin{itemize}
\raggedright
  \item \emph{Say it:} \emph{resonō, resonāre}
  \item \emph{Say it:} \emph{resonō, resonāre}, once more
  \item \emph{Say it:} a sentence of your own with \emph{resonō, resonāre}
  \item \emph{Recall:} say the Latin for a sound, then the Latin for to sound, then say \emph{resonō, resonāre} again
\end{itemize}

\begin{tcolorbox}[breakable,colback=teal!4,colframe=teal!35!black,title={Before you move on}]
What does resonō, resonāre mean? (\textquotedblleft{}To resound, echo\textquotedblright{}.) Say it once more.
\end{tcolorbox}
\section[symphōnia, symphōniae]{symphōnia, symphōniae — a band of musicians; harmony}

\label{lesson:LA-C545-symphonia}



\begin{quote}
Before the new one: say the Latin for to sound, then the Latin for to resound.
\end{quote}

\subsection*{You'll want to know: symphōnia, symphōniae}
\textbf{symphōnia, symphōniae} — \textquotedblleft{}a band of musicians; harmony\textquotedblright{}.

\begin{quote}
In convīviō symphōnia canit.
\end{quote}

\begin{quote}
\emph{At the dinner party a band plays.}
\end{quote}

Feminine, first declension, from Greek: sounds together, then the slaves who made music at dinner. English \emph{symphony}.

\begin{cousinweb}
\textbf{symphōnia} was borrowed from Greek \emph{symphonia}, built on \emph{syn} \textquotedblleft{}with\textquotedblright{} and \emph{phone} \textquotedblleft{}sound\textquotedblright{} (EtymDB; Lewis \& Short). English \emph{symphony} comes from it.
\end{cousinweb}

\begin{grammarlens}[title={What you've built}]
One more word for telling, explaining and describing at B1.
\end{grammarlens}

\subsection*{Your turn}
\begin{itemize}
\raggedright
  \item \emph{Say it:} \emph{symphōnia, symphōniae}
  \item \emph{Say it:} \emph{symphōnia, symphōniae}, once more
  \item \emph{Say it:} a sentence of your own with \emph{symphōnia, symphōniae}
  \item \emph{Recall:} say the Latin for to sound, then the Latin for to resound, then say \emph{symphōnia, symphōniae} again
\end{itemize}

\begin{tcolorbox}[breakable,colback=teal!4,colframe=teal!35!black,title={Before you move on}]
What does symphōnia, symphōniae mean? (\textquotedblleft{}A band of musicians; harmony\textquotedblright{}.) Say it once more.
\end{tcolorbox}
\section[concentus, concentūs]{concentus, concentūs — harmony, singing together}

\label{lesson:LA-C545-concentus}



\begin{quote}
Before the new one: say the Latin for to resound, then the Latin for a band of musicians.
\end{quote}

\subsection*{You'll want to know: concentus, concentūs}
\textbf{concentus, concentūs} — \textquotedblleft{}harmony, singing together\textquotedblright{}.

\begin{quote}
Concentus avium in silvā dulcis est.
\end{quote}

\begin{quote}
\emph{The chorus of the birds in the wood is sweet.}
\end{quote}

Masculine, fourth declension, from \emph{concinō}. Also agreement between people.

\begin{cousinweb}
\textbf{concentus} is the action noun of \emph{concinō} \textquotedblleft{}to sing together\textquotedblright{} (EtymDB; Lewis \& Short).
\end{cousinweb}

\begin{grammarlens}[title={What you've built}]
One more word for telling, explaining and describing at B1.
\end{grammarlens}

\subsection*{Your turn}
\begin{itemize}
\raggedright
  \item \emph{Say it:} \emph{concentus, concentūs}
  \item \emph{Say it:} \emph{concentus, concentūs}, once more
  \item \emph{Say it:} a sentence of your own with \emph{concentus, concentūs}
  \item \emph{Recall:} say the Latin for to resound, then the Latin for a band of musicians, then say \emph{concentus, concentūs} again
\end{itemize}

\begin{tcolorbox}[breakable,colback=teal!4,colframe=teal!35!black,title={Before you move on}]
What does concentus, concentūs mean? (\textquotedblleft{}Harmony, singing together\textquotedblright{}.) Say it once more.
\end{tcolorbox}
